%%%%%%%%%%%%%%%%%%%%%%%%%%%%%%%%%%%%%%%%%%%%%%%%%%%%%%%%%%
%%
%% MARK: Title
%%
%%%%%%%%%%%%%%%%%%%%%%%%%%%%%%%%%%%%%%%%%%%%%%%%%%%%%%%%%%

\chapter*{\S B31. Foliations and Diffeology}
\setcounter{footnote}{0}

%%%%%%%%%%%%%%%%%%%%%%%%%%%%%%%%%%%%%%%%%%%%%%%%%%%%%%%%%%
%%
%% MARK: Abstract
%%
%%%%%%%%%%%%%%%%%%%%%%%%%%%%%%%%%%%%%%%%%%%%%%%%%%%%%%%%%%

\begin{abstract}
  Spaces of leaves of foliations on manifolds are among the first examples we can think of when it comes to diffeology.
  Paradoxically,
  with the exception of Kronecker's foliation of a torus by an irrational line,
  I have not delved into this class of examples in the textbook.
  In \cite{IL90} we treated the particular case of the foliation of an $n$-dimensional torus $\Torus^n$ by an irrational hyperplane.
  In this note we go a little further into what diffeology can say about the space of leaves of a foliated manifold.
\end{abstract}

%%%%%%%%%%%%%%%%%%%%%%%%%%%%%%%%%%%%%%%%%%%%%%%%%%%%%%%%%%
%%
%% MARK: Content
%%
%%%%%%%%%%%%%%%%%%%%%%%%%%%%%%%%%%%%%%%%%%%%%%%%%%%%%%%%%%

Let us begin by recalling what a foliation is \cite{Law74}.
Let $M$ be an $n$-manifold \cite[4.1]{PIZ13}.
We say that a partition $\cL$ of $M$ into connected subspaces
is a {\em foliation} if there exists an atlas of charts
$$
  \pphi : \RR^k \times \RR^m \supset \U \to M, \quad \text{with } k = n-m,
$$
such that for every $L \in \cL$,
the preimage under $\pphi$ of each connected component of $L \cap \pphi(\U)$ is a vertical slice $\U_r = \U \cap [\{r\} \times \RR^m]$,
for some $r \in \RR^k$,
see \cite[1.33, 5.7]{PIZ13}.
In simple words,
a foliation is a partition that locally looks like a product of Euclidean spaces.

The elements $L$ of $\cL$ are called the {\em leaves} of the foliation.
The charts $\pphi$ above are called {\em adapted} to the foliation $\cL$.
Note that the leaves,
equipped with the subset diffeology,
are $m$-dimensional submanifolds \cite[4.4]{PIZ13}.
Indeed,
with the notations above,
the restrictions $\pphi \restriction \U_r$ are local diffeomorphisms
from $\RR^m$ to $L$.

\begin{Article}{Transversals.}{}
  We call a {\em transversal} of $\cL$ any $k$-plot $Q : V \to M$ such that there exists an adapted chart $\pphi : V \to M$ for which $Q(r) = \pphi(r,0)$.
  A plot $P$ in $M$ is said to be {\em transversal} if it takes its values in the image of a transversal $Q$.
\end{Article}

\begin{figure}[htb]
  \includegraphics[width=0.7\textwidth]{Figures/foliation.jpg}
  %\vspace{-1.0\baselineskip}\n
  \caption{Rectification of the flow.}
  \label{Rectification of the flow}
\end{figure}

Now,
let $\cL$ be equipped with the quotient diffeology \cite[1.50]{PIZ13},
and let $\pi : M \to \cL$ be the projection,
that is,
$\pi(x) = L$ if and only if $x \in L \in \cL$.

\begin{Article}{The space of leaves.}{}
  Let $M$ be a second countable Hausdorff manifold,
  and let $\cL$ be a foliation on $M$.
  \begin{enumerate}
    \item Let $Q : V \to M$ be a transversal.
      Then for every $L \in \cL$,
      the set $L \cap Q(V)$ (or equivalently $\{r \in V \mid Q(r) \in L\}$)
      is at most countable,
      hence either empty or (diffeologically) discrete.
    \item There exist generating families of $\cL$,
      equipped with the quotient diffeology,
      consisting of parametrizations $\cQ = \pi \circ Q$,
      where $Q$ is a transversal plot.
    \item The dimension \cite[2.22]{PIZ13} of the quotient $\cL$ is constant and equal to $k = n-m$.
  \end{enumerate}
\end{Article}

\begin{proof}
  Since $Q$ is the restriction of an adapted chart $\pphi$,
  we have $Q(r) = \pphi(r,0)$.
  A leaf $L$ is locally given by fixing the first $k$ coordinates.
  The set $\{r \in V \mid Q(r) \in L\}$ consists of those $r$ for which the leaf through $\pphi(r,0)$ is exactly $L$.
  Because leaves are transverse to the transversal,
  this set is discrete in $V$,
  hence at most countable.
  The other statements follow directly from the definition of the quotient diffeology and the local product structure.
\end{proof}
